\documentclass[11pt]{amsart}

\usepackage[T1]{fontenc}
\usepackage{lmodern}
\usepackage{microtype}
\usepackage[margin=1.05in]{geometry}
\usepackage{amsmath,amssymb,mathtools}
\usepackage{booktabs}
\usepackage{enumitem}
\usepackage{aliascnt}
\usepackage{graphicx}
\usepackage{tikz}
\usetikzlibrary{arrows.meta,calc}
\usepackage[numbers,sort&compress]{natbib}

\allowdisplaybreaks
\setlist{nosep,leftmargin=*}
\newtheorem{theorem}{Theorem}[section]
\newaliascnt{proposition}{theorem}
\newtheorem{proposition}[proposition]{Proposition}
\aliascntresetthe{proposition}
\newaliascnt{lemma}{theorem}
\newtheorem{lemma}[lemma]{Lemma}
\aliascntresetthe{lemma}
\newaliascnt{corollary}{theorem}
\newtheorem{corollary}[corollary]{Corollary}
\aliascntresetthe{corollary}
\theoremstyle{definition}
\newaliascnt{definition}{theorem}
\newtheorem{definition}[definition]{Definition}
\aliascntresetthe{definition}
\newaliascnt{remark}{theorem}
\newtheorem{remark}[remark]{Remark}
\aliascntresetthe{remark}

\usepackage[hidelinks]{hyperref}
\hypersetup{
  pdftitle={The Exact Region Formula for Arrangements of Constrained Long-Legged A's},
  pdfauthor={Siddhartha Mahajan and Paras Chopra},
  pdfsubject={Exact extremal region count for constrained long-legged A arrangements},
  pdfkeywords={plane arrangements, extremal graph theory, Mantel's theorem, OEIS A397182}
}
\usepackage[nameinlink,noabbrev]{cleveref}
\crefname{theorem}{theorem}{theorems}
\crefname{proposition}{proposition}{propositions}
\crefname{lemma}{lemma}{lemmas}
\crefname{corollary}{corollary}{corollaries}
\crefname{remark}{remark}{remarks}

\newcommand{\R}{\mathbb R}
\newcommand{\Q}{\mathbb Q}
\newcommand{\evec}{\mathbf e}
\newcommand{\B}{\mathcal B}
\newcommand{\eps}{\varepsilon}
\DeclareMathOperator{\detm}{det}

\title[Constrained long-legged A arrangements]{The Exact Region Formula for Arrangements of Constrained Long-Legged A's}

\author{Siddhartha Mahajan}
\address{Research Fellow, Lossfunk}
\email{siddharthamahajan03@gmail.com}

\author{Paras Chopra}
\address{Lossfunk}
\email{paras@lossfunk.com}

\date{23 July 2026}
\subjclass[2020]{Primary 05C10; Secondary 52C30, 05C35}
\keywords{plane arrangements, extremal graph theory, Mantel's theorem, geometric construction, OEIS A397182}

\begin{document}

\begin{abstract}
A constrained long-legged A consists of two rays with a common apex and a
crossbar joining points at equal distance from that apex.  Let $a(n)$ be the
maximum number of regions into which $n$ such objects divide the plane.
Cutler, Karlsson, and Sloane established that two constrained A's meet at most
eight times and that three of them cannot be pairwise eight-crossing.  The
graph of eight-crossing pairs is therefore triangle-free, and Mantel's
theorem gives
\[
 a(n)\le 7\binom n2+\left\lfloor\frac{n^2}{4}\right\rfloor+2n+1.
\]
We construct matching arrangements for every $n$.  The construction splits
the objects into a narrow family and a wide nested fan.  Every cross-family
pair has eight proper intersections, while every same-family pair has at
least seven.  Three elementary blow-ups, together with one cubic Taylor
calculation for the narrow family, verify all required component incidences.
Because these inequalities are strict, a generic rational perturbation
preserves them.  Consequently, for every $n\ge0$,
\[
 \boxed{\displaystyle
 a(n)=7\binom n2+\left\lfloor\frac{n^2}{4}\right\rfloor+2n+1
     =\left\lfloor\frac{15n^2-6n+4}{4}\right\rfloor.}
\]
\end{abstract}

\maketitle

\section{Introduction}

A constrained long-legged A is the union of two rays issuing from one apex
and a crossbar whose endpoints lie at equal distance from the apex.  The
region-maximization problem for these objects was introduced by Cutler,
Karlsson, and Sloane in their study of planar cuts by exotic knives
\cite{CKS2026}; the resulting sequence is OEIS A397182
\cite{OEISA397182}.  The equality of the two leg lengths is a metric
condition, so affine deformations that trivialize many line-arrangement
problems are unavailable here.

Let $a(n)$ denote the largest number of connected components of the
complement of a union of $n$ constrained A's.  The first values found in
\cite{CKS2026} are
\[
 a(0),a(1),a(2),a(3)=1,3,13,30.
\]
The main result of this paper determines the sequence completely.

\begin{theorem}\label{thm:main}
For every integer $n\ge0$,
\begin{equation}\label{eq:mainformula}
 a(n)=7\binom n2+\left\lfloor\frac{n^2}{4}\right\rfloor+2n+1
 =\left\lfloor\frac{15n^2-6n+4}{4}\right\rfloor.
\end{equation}
\end{theorem}

The upper bound is an extremal-graph argument.  Join two objects when they
meet eight times.  The three-object result of \cite{CKS2026} makes this
graph triangle-free, so Mantel's theorem bounds the number of eight-pairs.
Equality forces a balanced complete bipartite graph of eight-pairs and seven
crossings on every remaining pair.

The lower construction realizes exactly that colored graph.  For prescribed
part sizes $p,q$, we build $p$ narrow A's and $q$ wide A's.  The geometry is
organized around two points, $O=(0,0)$ and $C=(1,0)$, and is read on three
scales.  Near $O$, the narrow legs form a regular fan.  Near $C$, the narrow
crossbars become a short vertical stack and the high legs and crossbars of
the wide family become a nested fan.  Near the wide apices, a second blow-up
orders the low legs.  These limits establish seven incidences on every
same-family pair and eight on every cross-family pair.  The only limiting
incidences that land at component endpoints are two narrow bar--leg hits;
a third-order expansion places both strictly inside the bars.

The proof is entirely analytic.  Exact rational certificates through
$n=16$, which arose during discovery, are included as supplementary checks
but are not used in the proof of \cref{thm:main}.  Colored extremal graphs
also play a central role in the related published lollipop-bounds paper
\cite{MahajanChopra2026}.

\section{The geometric model and the upper bound}

\subsection{Constrained A's and proper crossings}

For an angle $\phi$, write
\[
 \evec(\phi)=(\cos\phi,\sin\phi).
\]
Given an apex $P\in\R^2$, a radius $r>0$, and two distinct unit vectors
$u,v$, define
\begin{equation}\label{eq:Adef}
 A(P,r;u,v)=
 \{P+tu:t\ge0\}\cup\{P+tv:t\ge0\}
 \cup[P+ru,P+rv].
\end{equation}
The first two components are the legs and the last is the crossbar.  We
always label the legs so that $\detm(u,v)>0$; their counterclockwise angular
separation then lies in $(0,\pi)$.

An intersection of components belonging to different A's is \emph{proper}
when it lies strictly past the apex on every participating leg and strictly
inside every participating crossbar.  An arrangement is \emph{generic} if
supporting lines of components from different A's are nonparallel, every
actual cross-shape component intersection is proper, no vertex lies on a
foreign component, and no point lies on components from three or more
A's.  It is enough to maximize over generic arrangements: transverse proper
crossings persist under perturbation, while a local generic resolution of a
parallelism, vertex incidence, or multiple crossing weakly increases the
number of complementary regions.  We will construct an open set with the
desired pairwise incidences and then choose a generic point in that set.

\subsection{The region identity}

\begin{proposition}\label{prop:regions}
A generic arrangement of $n$ constrained A's with $I$ finite crossings has
\[
 F=I+2n+1
\]
complementary regions.
\end{proposition}

\begin{proof}
Compactify the plane by one point at infinity.  One constrained A becomes a
connected graph with four vertices (the apex, the two crossbar endpoints,
and infinity) and five edges, hence first Betti number two.  The union of
$n$ disjoint A's with their points at infinity identified has first Betti
number $2n$.  Each proper transverse crossing identifies two interior
points of distinct arcs and increases the first Betti number by one.  Thus
the compactified union has first Betti number $2n+I$.  Euler's formula on
the sphere gives $F=(2n+I)+1$.
\end{proof}

\subsection{The eight-crossing graph}

We use the following two geometric facts established in \cite[Sec.~13.2]{CKS2026}.

\begin{proposition}[Cutler--Karlsson--Sloane]\label{prop:localfacts}
Two constrained long-legged A's have at most eight proper finite
intersections.  Moreover, no three constrained A's are pairwise
 eight-crossing.
\end{proposition}

For a generic arrangement, let $G_8$ be the graph whose vertices are the A's
and whose edges are the pairs with eight crossings.  By
\cref{prop:localfacts}, $G_8$ is triangle-free.

\begin{theorem}[Universal upper bound]\label{thm:upper}
For every $n\ge0$,
\begin{equation}\label{eq:upper}
 a(n)\le 7\binom n2+\left\lfloor\frac{n^2}{4}\right\rfloor+2n+1.
\end{equation}
If equality holds, the eight-crossing graph is a balanced complete bipartite
graph and every nonedge corresponds to a seven-crossing pair.
\end{theorem}

\begin{proof}
Let $e=e(G_8)$.  Every edge contributes at most eight crossings and every
nonedge at most seven, so
\[
 I\le 8e+7\left(\binom n2-e\right)=7\binom n2+e.
\]
Mantel's theorem \cite{Mantel1907} gives
$e\le\lfloor n^2/4\rfloor$.  Combining this inequality with
\cref{prop:regions} proves \eqref{eq:upper}.  Equality requires equality in
every step.  The equality case of Mantel's theorem is the balanced complete
bipartite graph, and every nonedge must contribute seven crossings.
\end{proof}

It remains to realize the equality pattern for arbitrary part sizes.

\section{A stability principle for component incidences}

We will repeatedly inspect a family of components after an affine rescaling
that depends on a small parameter $\eps>0$.  The following elementary
observation formalizes the limiting arguments.

\begin{lemma}[Blow-up stability]\label{lem:stability}
Let two supporting lines and their directed component domains depend
continuously on $\eps>0$.  Suppose that after an invertible affine change of
coordinates they converge, as $\eps\to0^+$, to two nonparallel lines whose
intersection lies in the relative interior of both limiting domains.  Then,
for all sufficiently small $\eps>0$, the original components have a proper
intersection.
\end{lemma}

\begin{proof}
The intersection parameters of two nonparallel lines are continuous
rational functions of their origins and directions.  At the limiting
configuration both parameters satisfy strict domain inequalities.  Those
inequalities therefore persist for all sufficiently small $\eps$.
\end{proof}

All asymptotics below are taken with the finite integers $p,q$ fixed.  Thus
an $O(\eps^r)$ estimate is uniform over every index occurring in the
construction, and finitely many smallness requirements can be met by one
choice of $\eps$.

For two A's, order the components as (first leg, second leg, crossbar) and
write a $3\times3$ $0$--$1$ matrix for their proper component incidences.
Our construction will establish the following entrywise lower bounds:
\begin{equation}\label{eq:targetmasks}
 M_{NN}=\begin{pmatrix}1&1&0\\1&1&0\\1&1&1\end{pmatrix},\qquad
 M_{NW}=\begin{pmatrix}1&1&1\\1&1&1\\0&1&1\end{pmatrix},\qquad
 M_{WW}=\begin{pmatrix}1&1&1\\0&1&1\\0&1&1\end{pmatrix}.
\end{equation}
They contain respectively seven, eight, and seven required incidences.  A
zero in \eqref{eq:targetmasks} means only that no incidence is required
there.

\section{The two-family construction}

Fix integers $p,q\ge1$ and put $m=\max\{p,q\}$.  Let
\[
 O=(0,0),\qquad C=(1,0),\qquad h=\frac32.
\]
We define $p$ narrow A's $N_1,\ldots,N_p$ and $q$ wide A's
$W_1,\ldots,W_q$, depending on a sufficiently small parameter $\eps>0$.

\subsection{The narrow family}

For $1\le i\le p$, set
\[
 a_i=3i,
 \qquad \alpha_i=a_i\eps,
 \qquad r_i=\cos\alpha_i,
 \qquad P_i=C-r_i\evec(\alpha_i).
\]
The two leg directions of $N_i$ are
\[
 d_i^- =\evec(\alpha_i-\eps),
 \qquad
 d_i^+ =\evec(\alpha_i+\eps),
\]
and the crossbar joins $P_i+r_i d_i^-$ to $P_i+r_i d_i^+$.
For small $\eps$, $r_i>0$ and this is a constrained A.

It is useful to write its crossbar in centered form.  Define
\[
 n_i=\evec(\alpha_i+\tfrac\pi2),\qquad
 M_i=C-r_i(1-\cos\eps)\evec(\alpha_i),\qquad
 s_i=r_i\sin\eps.
\]
Then the crossbar is
\begin{equation}\label{eq:narrowbar}
 \B_i=\{M_i+t n_i:-s_i\le t\le s_i\}.
\end{equation}
Its supporting line has the particularly simple equation
\begin{equation}\label{eq:narrowbarline}
 \evec(\alpha_i)\mathbin\cdot X=\cos\alpha_i\cos\eps.
\end{equation}

\subsection{The wide family}

For $1\le j\le q$, put
\begin{equation}\label{eq:wideparameters}
 q_j=4^{-j},\qquad
 \theta_j=2\arctan q_j,\qquad
 D_j=q_j^{-1}=4^j,
 \qquad L_j=4m+j,
 \qquad \lambda_j=L_j\eps.
\end{equation}
The low and high leg directions are
\begin{equation}\label{eq:widedirections}
 \ell_j=\evec(\lambda_j),
 \qquad
 u_j=\evec(\pi-\theta_j+\lambda_j).
\end{equation}
Place the high crossbar endpoint at
\[
 H_j=C+\eps(-D_j,h)
\]
and define
\begin{equation}\label{eq:wideradius}
 R_j=\frac{1+h\eps}{\sin(\theta_j-\lambda_j)},
 \qquad
 Q_j=H_j-R_j u_j.
\end{equation}
Then $(Q_j)_y=-1$.  The wide A $W_j$ has apex $Q_j$, common leg length
$R_j$, low leg direction $\ell_j$, and high leg direction $u_j$.  Its low
crossbar endpoint is
\[
 K_j=Q_j+R_j\ell_j.
\]
For small $\eps$, $0<\lambda_j<\theta_j$, so $R_j>0$ and the two directions
in \eqref{eq:widedirections} are distinct with positive orientation.

At $\eps=0$ the relevant apex and endpoint abscissae are
\begin{equation}\label{eq:AB}
 A_j=1+\cot\theta_j
     =1+\frac{D_j-q_j}{2},
 \qquad
 B_j=1+\cot\frac{\theta_j}{2}=1+D_j.
\end{equation}
Indeed,
\begin{equation}\label{eq:widelimits}
 Q_j\longrightarrow(A_j,-1),\qquad
 K_j\longrightarrow(B_j,-1),\qquad
 H_j\longrightarrow C.
\end{equation}

\begin{figure}[t]
\centering
\resizebox{\textwidth}{!}{%
\begin{tikzpicture}[x=0.88cm,y=0.88cm,>=Latex,font=\small]
  % left panel: C blow-up
  \begin{scope}[shift={(0,0)}]
    \draw[->] (-4.6,0)--(4.3,0) node[right] {$X$};
    \draw[->] (0,-2.7)--(0,3.0) node[above] {$Y$};
    \draw[very thick] (0,-1)--(0,1);
    \draw[densely dashed] (-4.4,1)--(4.0,1) node[right] {$N_i^+$};
    \draw[densely dashed] (-4.4,-1)--(4.0,-1) node[right] {$N_i^-$};
    \coordinate (H1) at (-1.1,1.5);
    \coordinate (H2) at (-2.5,1.5);
    \fill (H1) circle (1.2pt) node[above] {$(-D_j,h)$};
    \fill (H2) circle (1.2pt);
    \draw[thick] ($(H1)+(-2.0,1.2)$)--($(H1)+(4.0,-2.4)$);
    \draw[thick] ($(H1)+(-2.0,0.55)$)--($(H1)+(4.0,-1.1)$);
    \draw[thick] ($(H2)+(-1.4,0.45)$)--($(H2)+(5.0,-1.6)$);
    \draw[thick] ($(H2)+(-1.4,0.20)$)--($(H2)+(5.0,-0.72)$);
    \node[align=center] at (0,-3.35) {(a) blow-up at $C$: narrow bar/legs\\and the nested high/bar fan};
  \end{scope}
  % right panel: y=-1 blow-up
  \begin{scope}[shift={(11.4,0)}]
    \draw[->] (-0.3,0)--(7.6,0) node[right] {$x$};
    \draw[->] (0,-0.3)--(0,3.0) node[above] {$Y$};
    \draw[thick] (0.7,0)--(6.5,2.7);
    \draw[thick] (2.2,0)--(6.6,2.7);
    \draw[very thick] (2.2,0)--(2.2,2.1);
    \draw[very thick] (4.8,0)--(4.8,2.5);
    \node[below] at (0.7,0) {$A_j$};
    \node[below] at (2.2,0) {$A_k$};
    \node[below] at (4.8,0) {$B_k$};
    \node[align=center] at (3.7,-3.35) {(b) blow-up at $y=-1$: low legs\\meet the next apex and low endpoint};
  \end{scope}
\end{tikzpicture}%
}
\caption{The two limiting charts used for narrow--wide and wide--wide
incidences.  The drawing is schematic; only the strict order and slope
relations are used.}
\label{fig:blowups}
\end{figure}

\section{Verification of the three pair types}

\subsection{Narrow--narrow pairs}

\begin{lemma}\label{lem:NN}
For $1\le i<k\le p$ and all sufficiently small $\eps>0$, the component
incidence matrix of $(N_i,N_k)$ contains $M_{NN}$ from
\eqref{eq:targetmasks}.  In particular, the pair has at least seven proper
intersections.
\end{lemma}

\begin{proof}
Write $a=a_i$, $b=a_k$, so $b-a\ge3$.
First apply the blow-up $(x,y)\mapsto(x,y/\eps)$ near $O$.  The leg of $N_i$
with sign $\sigma\in\{-1,+1\}$ converges to the directed line
\begin{equation}\label{eq:narrowleglimit}
 Y=-a+(a+\sigma)x,\qquad x>0.
\end{equation}
For signs $\sigma,\tau$, the limiting intersection of the corresponding
legs of $N_i,N_k$ has
\[
 x=\frac{b-a}{b-a+\tau-\sigma}>0,
\]
because the denominator is at least $(b-a)-2\ge1$.  By
\cref{lem:stability}, all four leg--leg intersections are proper.

Next consider the intersection of $\B_i$ with the $\tau$-leg of $N_k$.  Put
$A=a\eps$, $B=b\eps$ and write the point on the bar as $M_i+t n_i$.  Direct
line intersection gives a ray parameter
\begin{equation}\label{eq:ttau}
 \rho_\tau=
 \frac{\cos A(\cos\eps-1)+\cos B\cos(B-A)}
      {\cos(B-A+\tau\eps)}
 =1+O(\eps^2)
\end{equation}
and bar coordinate
\begin{equation}\label{eq:qtau}
 t_\tau=-\cos B\sin(B-A)
       +\rho_\tau\sin(B-A+\tau\eps).
\end{equation}
Taylor expansion at $\eps=0$ yields
\begin{align}
 t_{-}+s_i
   &=\frac{(b-a)(2a+1)}{2}\eps^3+O(\eps^5),
 &s_i-t_{-}&=2\eps+O(\eps^3),\label{eq:minusbar}\\
 s_i-t_{+}
   &=\frac{(b-a)(2a-1)}{2}\eps^3+O(\eps^5),
 &t_{+}+s_i&=2\eps+O(\eps^3).\label{eq:plusbar}
\end{align}
Here $a=3i\ge3$, so all four displayed quantities are positive for small
$\eps$.  Together with \eqref{eq:ttau}, this proves that $\B_i$ meets both
legs of $N_k$ properly.

Finally, by \eqref{eq:narrowbarline}, every narrow crossbar supporting line
passes through
\[
 X_\eps=(\cos\eps,0).
\]
Its coordinate on $\B_i$ is
\[
 n_i\mathbin\cdot(X_\eps-M_i)
 =\sin(a\eps)(1-\cos\eps)=O(\eps^3),
\]
whereas $s_i=\eps+O(\eps^3)$.  Hence $X_\eps$ lies strictly inside both
$\B_i$ and $\B_k$ for small $\eps$.  We have obtained the seven entries of
$M_{NN}$.
\end{proof}

\begin{remark}\label{rem:concurrence}
The unperturbed narrow crossbars are concurrent at $X_\eps$.  This is a
convenient feature of the proof, not a claimed general-position
configuration.  The concurrence will be removed in \cref{sec:perturbation}
without changing any of the strict incidences above.
\end{remark}

\subsection{Narrow--wide pairs}

\begin{lemma}\label{lem:NW}
For $1\le i\le p$ and $1\le j\le q$, and all sufficiently small
$\eps>0$, the component incidence matrix of $(N_i,W_j)$ contains $M_{NW}$
from \eqref{eq:targetmasks}.  Thus the pair has eight proper intersections.
\end{lemma}

\begin{proof}
Set $a=a_i$, $L=L_j$, $D=D_j$, and abbreviate
\begin{equation}\label{eq:muv}
 \mu=\tan\theta_j=\frac{2q_j}{1-q_j^2},
 \qquad
 \nu=\tan\frac{\theta_j}{2}=q_j.
\end{equation}
We first meet the wide low leg.  The determinant of its direction with the
$\sigma$-leg of $N_i$ is
\[
 \sin\bigl((L-a-\sigma)\eps\bigr),
\]
and
\[
 L-a-\sigma\ge(4m+1)-(3m+1)=m>0.
\]
If $r_\sigma$ and $s_\sigma$ are the two ray parameters at their
intersection, the determinant formula for two lines gives
\begin{equation}\label{eq:farparameters}
 \eps r_\sigma\longrightarrow\frac1{L-a-\sigma},
 \qquad
 \eps s_\sigma\longrightarrow\frac1{L-a-\sigma}.
\end{equation}
Both are positive, so each narrow leg meets the wide low leg properly.

The remaining six required incidences are visible in the blow-up at $C$,
\[
 (X,Y)=\left(\frac{x-1}{\eps},\frac y\eps\right).
\]
The two narrow legs tend to $Y=-1$ and $Y=1$, while $\B_i$ tends to the open
vertical segment $X=0$, $-1<Y<1$.  The high leg and crossbar of $W_j$ tend
to
\begin{align}
 \mathcal H_j&:Y-h=-\mu(X+D),\label{eq:highlimit}\\
 \mathcal C_j&:Y-h=-\nu(X+D),\qquad X>-D,\label{eq:barlimit}
\end{align}
respectively.  The inequality in \eqref{eq:barlimit} is the limiting
interior direction from the high endpoint toward the low endpoint.

For a narrow leg $Y=\sigma$, its intersection with \eqref{eq:barlimit}
satisfies
\[
 X+D=\frac{h-\sigma}{\nu}>0
\]
because $h=3/2>1$.  Hence both narrow legs meet the wide bar in its
interior.  In this chart the apex of the high ray escapes to infinity along
$\mathcal H_j$ in the direction opposite the ray.  Consequently every fixed
finite point of $\mathcal H_j$, including its intersections with the two
narrow legs, lies in the ray's interior for small $\eps$.

At $X=0$, the limiting heights of the narrow-bar intersections are
\begin{equation}\label{eq:narrowwideheights}
 h-D\mu=h-\frac{2}{1-q_j^2},
 \qquad
 h-D\nu=h-1=\frac12.
\end{equation}
Since $0<q_j\le1/4$,
\[
 -\frac{19}{30}
 =\frac32-\frac{32}{15}
 \le h-D\mu< -\frac12,
\]
so both values in \eqref{eq:narrowwideheights} lie strictly between $-1$
and $1$.  The second intersection also has $X+D=D>0$, placing it inside the
wide bar.  By \cref{lem:stability}, the narrow bar meets both the wide high
leg and the wide bar properly.  Together with \eqref{eq:farparameters},
these are exactly the eight required entries of $M_{NW}$.
\end{proof}

\subsection{Wide--wide pairs}

\begin{lemma}\label{lem:WW}
For $1\le j<k\le q$ and all sufficiently small $\eps>0$, the component
incidence matrix of $(W_j,W_k)$ contains $M_{WW}$ from
\eqref{eq:targetmasks}.  In particular, the pair has at least seven proper
intersections.
\end{lemma}

\begin{proof}
Write
\[
 \mu_r=\tan\theta_r=\frac{2q_r}{1-q_r^2},
 \qquad \nu_r=q_r
 \qquad(r=j,k).
\]
Because $q_k\le q_j/4$ and $q_j\le1/4$,
\begin{equation}\label{eq:slopechain}
 \mu_j>\nu_j>\mu_k>\nu_k>0.
\end{equation}
Indeed, only the middle inequality needs comment, and
\[
 \mu_k\le\frac{q_j/2}{1-q_j^2/16}<q_j=\nu_j.
\]

In the blow-up at $C$, each high leg or crossbar of $W_r$ has a limiting
line through $(-D_r,h)$ of slope $-s$, where $s$ is respectively $\mu_r$ or
$\nu_r$.  Consider one such line for $r=j$ with slope $-s$ and one for
$r=k$ with slope $-t$.  By \eqref{eq:slopechain}, $s>t$.  If their
intersection has displacements
\[
 U_j=X+D_j,
 \qquad U_k=X+D_k,
\]
then, with $\Delta=D_k-D_j>0$,
\begin{equation}\label{eq:fanintersection}
 U_j=\frac{t\Delta}{s-t}>0,
 \qquad
 U_k=\frac{s\Delta}{s-t}>0.
\end{equation}
Whenever one of the two components is a crossbar, the corresponding
positive displacement places the intersection strictly beyond its high
endpoint toward its low endpoint.  The displacement is finite in the
blow-up, whereas the other endpoint escapes to infinity, so it is strictly
inside the bar.  The apex of a high leg escapes to infinity along its limiting line in
the direction opposite the ray, so every fixed finite point on that line is
in the ray's interior for small $\eps$.  Thus
\eqref{eq:fanintersection} proves the four high/bar incidences
\[
 (\text{high},\text{high}),\quad
 (\text{high},\text{bar}),\quad
 (\text{bar},\text{high}),\quad
 (\text{bar},\text{bar}).
\]

It remains to show that the low leg of $W_j$ meets all three components of
$W_k$.  Use the second blow-up
\[
 (x,Y)=\left(x,\frac{y+1}{\eps}\right).
\]
By \eqref{eq:AB}--\eqref{eq:widelimits}, the low leg of $W_r$ tends to
\begin{equation}\label{eq:lowlimit}
 Y=L_r(x-A_r),\qquad x>A_r.
\end{equation}
The portion of the high ray of $W_k$ at height $O(\eps)$ tends to the
vertical half-line
\[
 x=A_k,\qquad Y>0.
\]
The low endpoint of the $k$th crossbar has the more precise limit
\begin{equation}\label{eq:lowendpointheight}
 \left(B_k,\kappa_k\right),
 \qquad
 \kappa_k=\frac{L_k}{\sin\theta_k}
          =\frac{L_k(D_k+q_k)}{2},
\end{equation}
because $((K_k)_y+1)/\eps=R_k\sin\lambda_k/\eps\to
L_k/\sin\theta_k$.  The portion of the crossbar near that endpoint tends
to the vertical half-line $x=B_k$, $Y>\kappa_k$.

Since $j<k$, we have $A_j<A_k<B_k$ and $L_j<L_k$.  The limiting
intersections of the $j$th low leg with the $k$th high leg and crossbar are
respectively
\begin{equation}\label{eq:lowverticalhits}
 (A_k,L_j(A_k-A_j)),
 \qquad
 (B_k,L_j(B_k-A_j)).
\end{equation}
The first has positive height and is forward of the $j$th apex.  For the
second, note that $k\ge2$, $D_j\le D_k/4$, $q_k^2\le1/256$, and
\[
 \frac{L_k}{L_j}
 =1+\frac{k-j}{4m+j}<\frac54.
\]
Therefore
\begin{align*}
 B_k-A_j
 &=D_k-\frac{D_j}{2}+\frac{q_j}{2}
   \ge\frac{7D_k}{8},\\
 \frac{\kappa_k}{L_j}
 &=\frac{L_k}{L_j}\,\frac{D_k+q_k}{2}
 <\frac54\cdot\frac{257}{256}\cdot\frac{D_k}{2}
 <\frac{7D_k}{8}.
\end{align*}
Thus $L_j(B_k-A_j)>\kappa_k$, so the second point in
\eqref{eq:lowverticalhits} lies strictly inside the limiting crossbar.
Both incidences are therefore proper for small $\eps$.

Finally, the two limiting low lines in \eqref{eq:lowlimit} meet at
\begin{equation}\label{eq:lowlowx}
 x_*=
 \frac{L_kA_k-L_jA_j}{L_k-L_j}
 =A_k+\frac{L_j(A_k-A_j)}{L_k-L_j}>A_k>A_j.
\end{equation}
Thus the low--low intersection lies forward on both rays.  These three low
incidences and the four incidences from \eqref{eq:fanintersection} give
$M_{WW}$.
\end{proof}

\section{Generic rational perturbation and completion of the proof}
\label{sec:perturbation}

The preceding construction was chosen for transparent limiting geometry,
not for general position.  In particular, the narrow crossbars concur as
noted in \cref{rem:concurrence}.  We now remove every degeneracy at once.

\begin{lemma}[Generic rational realization]\label{lem:generic}
For every $p,q\ge1$, there is a generic arrangement of $p+q$ constrained
A's, with rational apices, rational positive radii, and rational unit leg
directions, such that every narrow--wide pair has at least eight proper
intersections and every same-family pair has at least seven.
\end{lemma}

\begin{proof}
Choose $\eps>0$ small enough that all strict component-domain inequalities
proved in \cref{lem:NN,lem:NW,lem:WW} hold simultaneously.  They define an
open neighborhood $\mathcal U$ of the resulting point in the parameter
space of apices, radii, and leg directions.

Work in a local rational chart for each unit direction,
\[
 t\longmapsto
 \left(\frac{1-t^2}{1+t^2},\frac{2t}{1+t^2}\right).
\]
After denominators are cleared, each unwanted event---supporting-line
parallelism, a foreign component through an apex or crossbar endpoint, or
the coincidence of two distinctly labelled finite crossings---is contained
in the zero set of a polynomial in the apex coordinates, radii, and chart
parameters.  Only finitely many labels occur.  None of these polynomials is
identically zero on $\mathcal U$: an arbitrarily small independent change of
one participating apex, radius, or direction destroys the corresponding
event.  The bad locus is therefore a finite union of proper algebraic zero
sets and has empty interior.

Rational points are dense in the chart.  Since the bad locus is closed in
$\mathcal U$ and cannot contain a dense subset of $\mathcal U$, there is a
rational point of $\mathcal U$ outside it.  It gives a generic arrangement
with rational apices, positive rational radii, and rational unit directions.
All required incidences persist because their inequalities are strict.
\end{proof}

\begin{proof}[Proof of \cref{thm:main}]
The cases $n=0,1$ give $1$ and $3$ regions directly.  Let $n\ge2$ and take
\[
 p=\left\lfloor\frac n2\right\rfloor,
 \qquad
 q=\left\lceil\frac n2\right\rceil.
\]
Apply \cref{lem:generic}.  The $pq$ cross-family pairs contribute at least
eight crossings each, while the remaining pairs contribute at least seven.
Because the arrangement is generic, all these crossings are globally
distinct, and therefore
\begin{align*}
 I
 &\ge 8pq+7\left(\binom p2+\binom q2\right)\\
 &=7\binom n2+pq
 =7\binom n2+\left\lfloor\frac{n^2}{4}\right\rfloor.
\end{align*}
The reverse inequality is \cref{thm:upper}.  Hence equality holds, and
\cref{prop:regions} gives \eqref{eq:mainformula}.
\end{proof}

\begin{corollary}\label{cor:arbitraryparts}
For every $p,q\ge1$, there is a generic rational arrangement of $p+q$
constrained A's in which every cross-family pair has exactly eight crossings
and every same-family pair has exactly seven.  It has
\[
 7\binom{p+q}{2}+pq+2(p+q)+1
\]
regions.
\end{corollary}

\begin{proof}
The cross-family pairs have exactly eight crossings by
\cref{prop:localfacts,lem:generic}.  If a same-family pair also had eight,
then that pair together with any object from the other family would be three
pairwise eight-crossing A's, contrary to \cref{prop:localfacts}.  The region
count follows from \cref{prop:regions}.
\end{proof}

\begin{corollary}
The sequence begins
\[
\begin{split}
1,3,13,30,55,87,127,174,229,291,361,438,523,615,715,822,937,\ldots.
\end{split}
\]
For equality arrangements, the graph of eight-crossing pairs is necessarily
$K_{\lfloor n/2\rfloor,\lceil n/2\rceil}$ and every within-part pair has
exactly seven crossings.
\end{corollary}

\section{Exact computational certificates}
\label{sec:certificates}

The proof above does not depend on computation.  During discovery, however,
exact generic rational constructions were obtained separately for every
$4\le n\le16$.  They provide useful independent checks of the equality
pattern.  The supplementary archive contains the certificates listed in
\cref{tab:certificates} and a single verifier using only Python's
\texttt{fractions.Fraction} class.

For every certificate the verifier checks:
\begin{enumerate}[label=(\roman*)]
\item equal positive squared norms of the two leg directions and a positive
radius for every A;
\item all nine component-pair predicates for every pair of A's, with strict
ray and crossbar domains;
\item the balanced seven/eight pair-crossing matrix;
\item absence of cross-shape supporting-line parallelism;
\item absence of every foreign apex or crossbar-endpoint incidence; and
\item global distinctness of all proper rational crossings.
\end{enumerate}

\begin{table}[ht]
\centering
\caption{Exact rational certificates included in the supplement.}
\label{tab:certificates}
\begin{tabular}{@{}rrrrl@{}}
\toprule
$n$ & part sizes & crossings & regions & file\\
\midrule
4  & $2+2$ & 46  & 55  & \texttt{certificates/n04.json}\\
5  & $2+3$ & 76  & 87  & \texttt{certificates/n05.json}\\
6  & $3+3$ & 114 & 127 & \texttt{certificates/n06.json}\\
7  & $3+4$ & 159 & 174 & \texttt{certificates/n07.json}\\
8  & $4+4$ & 212 & 229 & \texttt{certificates/n08.json}\\
9  & $4+5$ & 272 & 291 & \texttt{certificates/n09.json}\\
10 & $5+5$ & 340 & 361 & \texttt{certificates/n10.json}\\
11 & $5+6$ & 415 & 438 & \texttt{certificates/n11.json}\\
12 & $6+6$ & 498 & 523 & \texttt{certificates/n12.json}\\
13 & $6+7$ & 588 & 615 & \texttt{certificates/n13.json}\\
14 & $7+7$ & 686 & 715 & \texttt{certificates/n14.json}\\
15 & $7+8$ & 791 & 822 & \texttt{certificates/n15.json}\\
16 & $8+8$ & 904 & 937 & \texttt{certificates/n16.json}\\
\bottomrule
\end{tabular}
\end{table}

A second standard-library script implements a fully rational, nongeneric
version of the two-family base construction.  It checks the exact masks
\[
 \texttt{110110111},\qquad
 \texttt{111111011},\qquad
 \texttt{111011011}
\]
for narrow--narrow, narrow--wide, and wide--wide pairs respectively.  The
archived run covers every balanced size $2\le n\le64$.  This is a regression
check rather than a logical input to the all-$n$ proof.

\section{Concluding remarks}

The upper and lower arguments isolate two different kinds of structure.
The upper bound remembers only whether a pair has seven or eight crossings;
its equality case is forced by Mantel's theorem.  The lower bound is more
geometric: the two color classes require different scales and different
local models.  The narrow family supplies four leg--leg crossings and three
bar incidences at a common macroscopic location.  The wide nested fan uses
an exponentially separated opening parameter $q_j=4^{-j}$ and a linearly
ordered drift $L_j\eps$.  The scale separation makes the three pair types
independent enough to verify by elementary limits.

More generally, \cref{cor:arbitraryparts} realizes every complete
bipartite graph $K_{p,q}$ as an eight-crossing graph, with seven crossings on
every remaining pair.  Thus the three masks in \eqref{eq:targetmasks} are
compatible at arbitrary scale.  The two blow-up centers separate the
relevant incidence constraints into two elementary line fans; the generic
perturbation then converts their strict local inequalities into a global
arrangement.

\bibliographystyle{amsplain}
\bibliography{references}

\end{document}
